\scanpage{104}
Variante globale : on définit des hypertor globaux \deleted{\unsure{par}} prenant $\mathbb R\Gamma_X\,L_{\bullet}\overset{\mathbb L}{\otimes}M_{\bullet}$.
% slanted
\margin{Mais il faut supposer, si on ne fait pas d'\unsure{hyp.} suppl. \unsure{pour|sur} $L_{\bullet},M_{\bullet}$, que $\Gamma_X$ est de dim \unsure{coh} finie !}

\uline{Lois formelles les plus évidentes} (lois internes).

1) Associativité de $\overset{\mathbb L}{\otimes}$

\unsure{2}) Commutativité (attention au signe !)

3)
% slanted
\margin{N.B. $\mathbb R\mathscr Hom(L_{\bullet},K^{\bullet})$ défini dès que $K^{\bullet}\in\mathrm{Ob}\,D^+(X)$, mais $L_{\bullet}$ peut être quelconque}
\diagram{\begin{tikzcd}[ampersand replacement=\&,column sep=.2em,row sep=.4em]
\mathbb R\mathscr Hom(L_{\bullet}\overset{L}{\otimes}M_{\bullet},K^{\bullet}) \& {}\simeq\mathbb R\mathscr Hom(L_{\bullet},\mathbb R\mathscr Hom(M_{\bullet},K^{\bullet}))\\
\& {}\simeq\mathbb R\mathscr Hom(M_{\bullet},\mathbb R\mathscr Hom(L_{\bullet},K^{\bullet}))\\
\mathbb R\mathrm{Hom}(L_{\bullet}\overset{\mathbb L}{\otimes}M_{\bullet},K^{\bullet})\arrow[uu,out=165,in=165] \& {}\simeq\mathbb R\mathrm{Hom}(L_{\bullet},\mathbb R\deleted{\unsure{H}}\mathscr Hom(M_{\bullet},K^{\bullet}))\\
\& {}\simeq\mathbb R\mathrm{Hom}(M_{\bullet},\mathbb R\mathscr Hom(L_{\bullet},K^{\bullet}))
\end{tikzcd}}
\[
\begin{aligned}
&\deleted{\mathrm{Hom}(L_{\bullet},\mathbb R\mathscr Hom}\\
\mathrm{Hom}(L_{\bullet}\overset{\mathbb L}{\otimes}M_{\bullet},K^{\bullet})&\simeq\mathrm{Hom}(L_{\bullet},\mathbb R\mathscr Hom(M_{\bullet},K^{\bullet})\\
&\simeq\mathrm{Hom}(M_{\bullet},\mathbb R\mathscr Hom(L_{\bullet},K^{\bullet}))
\end{aligned}
\]

\uline{Remarque sur changement de base non plat}

On définit $\mathbb L f^*:\deleted{\unsure{D}}D^-(Y)\longrightarrow D^-(X)$ encore en passant par $K^-(Y)^{\mathrm{plat}}/\text{quasi-isom}$, en notant que $f^*$ \deleted{\unsure{sur}} \deleted{\unsure{$K^-$}} envoie $K^-(Y)^{\unsure{\mathrm{plat}}}$ dans $K^-(X)^{\mathrm{plat}}$ et transforme qu.\ isom en idem. Avec cette définition, on aura encore
\[
\mathbb L f^*(L_{\bullet}\overset{\mathbb L}{\otimes}M_{\bullet})\simeq\mathbb L f^*(L_{\bullet})\overset{\mathbb L}{\otimes}\mathbb L f^*(M_{\bullet})
\]
Notons aussi formule de transitivité
% archive sheet number: 103
